%%%%%%%%%%%%%%%%%%%%%%%%%%%%%%%%%%%%%%%%%%%%%%%%%%%%%%%%%%%%%%%%%%%%%%%%%%%
%%%             TRABALHOS CORRELATOS               %%%
%%%%%%%%%%%%%%%%%%%%%%%%%%%%%%%%%%%%%%%%%%%%%%%%%%%%%%%%%%%%%%%%%%%%%%%%%%%

\setlength{\parskip}{0.3cm}

\chapter{Trabalhos Correlatos}
\label{ch:trabalhos_correlatos}
O presente capítulo tem como objetivo apresentar pesquisas que possuem relação
com a temática desta monografia. Esses trabalhos foram selecionados pela
convergência entre desenvolvimento do \gls{pc} e práticas gamificadas, com foco
no engajamento e no aprendizado de matemática. A análise desses trabalhos
oferece o embasamento teórico e prático necessário para o desenvolvimento da
solução proposta. O Quadro \ref{qua:correlatos} sintetiza a análise comparativa
dos trabalhos correlatos, categorizando cada estudo quanto aos conceitos de
\gls{pc}, gamificação e matemática, além de identificar as lacunas que o
presente trabalho propõe solucionar.

 
A dissertação de \cite{silva2020}, intitulada Pensamento computacional:
uma análise dos documentos oficiais e das questões de Matemática dos
vestibulares, consolida-se como estudo correlato para estabelecer o referencial
teórico ao investigar diretrizes educacionais, como a \gls{bncc}, buscando contribuir
com a discussão sobre o papel do \gls{pc} na Educação
Básica, servindo como base para o mapeamento dos conceitos do \gls{pc} envolvidos em questões
de matemática. Esse estudo não aborda aspectos sobre gamificação, sendo necessária 
a integração do tema.
  
O artigo de \cite{lais2023}, Enigma pirata: um jogo para exercitar Matemática e
Pensamento Computacional através da Teoria dos Conjuntos, descreve o
desenvolvimento de um jogo digital gamificado de aventura voltado para o ensino
e aprendizado da matemática e do \gls{pc}. A obra aborda o
processo de avaliação do engajamento, mas não descreve de forma completa o uso
da metodologia aplicada. 

O artigo de \cite{teodosio2023atividade} demonstra que a gamificação por meio de
planilhas eletrônicas e pixel art engaja os alunos através de estímulo visual e
descreve os pilares do \gls{pc} e elementos de \gls{pbl} empregados para engajamento e
resolução de problemas matemáticos. O estudo também relata o processo de
avaliação do engajamento, mas o artefato desenvolvido não é um jogo digital e
fica limitado ao software utilizado.

Embora focada na construção de artefatos digitais pelos próprios estudantes, a
pesquisa tem como objetivo validar como o \gls{pc} é utilizado para solucionar
problemas de matemática e física na construção dos jogos. \cite{filho2025}
destaca que os estudantes aplicaram conceitos do plano cartesiano para
posicionar os personagens e variáveis matemáticas para simular os efeitos de
aceleração e gravidade.


\begin{quadro}[H]
    \centering
    \caption{Análise comparativa dos trabalhos correlatos}
    \label{qua:correlatos}
    \renewcommand{\arraystretch}{1.2}
    \small
    \begin{tabular}{| >{\bfseries\raggedright\arraybackslash}p{2.3cm} |
        >{\raggedright\arraybackslash}p{2.3cm} |
        >{\raggedright\arraybackslash}p{3.5cm} |
        >{\raggedright\arraybackslash}p{3.0cm} |
        >{\raggedright\arraybackslash}p{3.0cm} |}
        \hline
        \normalfont\bfseries Trabalho & \normalfont\bfseries Conceito &
        \normalfont\bfseries Abordagem do \gls{pc} &
        \normalfont\bfseries Gamificação &
        \normalfont\bfseries Matemática \\
        \hline

        \cite{silva2020} &
        Mapeamento de conceitos de \gls{pc} em provas de matemática &
        Investiga como os pilares do \gls{pc} estão presentes em questões de vestibulares e nos documentos oficiais da \gls{bncc}. &
        \textbf{Não aborda.} &
        Análise de questões de vestibulares e diretrizes curriculares. \\
        \hline

        \cite{lais2023} &
        Jogo digital gamificado &
        Integra os pilares do \gls{pc} por meio da resolução de problemas de Teoria dos Conjuntos no contexto do jogo. &
        Jogo de aventura com elementos de progressão, mas sem detalhamento completo da metodologia de engajamento. &
        Exercita conteúdos matemáticos (Teoria dos Conjuntos) de forma contextualizada no enredo. \\
        \hline

        \cite{teodosio2023atividade} &
        Atividade gamificada &
        Desenvolve os pilares do \gls{pc} por meio de resolução de problemas matemáticos com apoio de planilhas eletrônicas e \textit{pixel art}. &
        Utiliza \gls{pbl} e estímulo visual, mas o artefato não é um jogo digital e fica limitado ao software de planilhas. &
        Aplica conceitos matemáticos em atividades gamificadas com enfoque em engajamento. \\
        \hline

        \cite{filho2025} &
        Construção de artefatos digitais &
        Valida a aplicação do \gls{pc} na resolução de problemas matemáticos e físicos durante a criação de jogos pelos estudantes. &
        Os próprios estudantes constroem elementos gamificados, mas não há mecânicas de recompensa ou progressão estruturadas. &
        Aborda matemática e física por meio do posicionamento de personagens e simulação de grandezas físicas. \\
        \hline

        Este trabalho &
        \textit{Software} gamificado educacional &
        Operacionaliza os quatro pilares do \gls{pc} (decomposição, abstração, reconhecimento de padrões e algoritmos) de forma sequencial e obrigatória na resolução de questões do \gls{enem}. &
        Implementa sistema completo de \gls{pbl} com \gls{xp}, insígnias, ligas, matriz de penalidades e custo estratégico de dicas. &
        Processa questões reais do \gls{enem} de matemática por meio de pipeline \gls{etl}, mapeando competências matemáticas aos pilares do \gls{pc}. \\
        \hline
    \end{tabular}

    \vspace{0.2cm}
    {\small Fonte: Elaborado pelo autor (2026).}
\end{quadro}

